\section{Ejercicios}

\subsection{Ejercicio 1: k caballos en tablero $n \times n$}
% TODO (5 pts):
% (a) variables de tu formulación CSP
% (b) dominio de cada variable
% (c) qué conjuntos de variables están restringidos y cómo
% (d) versión de "máximo número de caballos" como búsqueda local: Acciones,
%     Resultado, función objetivo

\subsection{Ejercicio 2: Sudoku con b\'usqueda local}
% TODO (5 pts): qué tan bien le iría a una heurística que minimiza
% restricciones violadas, sobre asignaciones completas.

\subsection{Ejercicio 3: \'arbol de juego del gato}
% TODO (5 pts): (a) número aproximado de partidas, (b) árbol completo hasta
% profundidad 2 con simetría, (c) evaluaciones en profundidad 2 con
% Eval(s) = 3*X2(s) + X1(s) - (3*O2(s) + O1(s)), (d) valores minimax
% propagados a profundidad 1 y 0 + mejor jugada inicial, (e) nodos podados
% por alfa-beta con orden óptimo.
% Sugerencia: una vez que tengas minimax/alpha_beta.py funcionando (sección
% de adversarios), correr el árbol real con tu propio código y contrastarlo
% contra lo que resuelvas a mano aquí es una buena forma de verificarte.

\subsection{Ejercicio 4: correctud de la poda alfa-beta}
% TODO (5 pts): (a)-(d), prueba formal siguiendo el enunciado.
